\documentclass[UTF8,11pt]{ctexart}
\usepackage[a4paper,margin=22mm]{geometry}
\usepackage{graphicx,enumitem,amssymb,tabularx}
\usepackage[normalem]{ulem}
\usepackage{xeCJKfntef}
\setlength{\parindent}{0pt}
\setlength{\parskip}{.6em}
\sloppy
\begin{document}
\section*{2024年12月英语四级真题(第1套)}
\subsection*{2024 年12 月大学英语四级考试真题（第1 套）}


\subsection*{\textbf{Part I Writing} \textbf{(30 minutes)}}


\textbf{Directions:} \textit{Suppose the business school of your university is conducting a survey to collect students’ opinions on} \textit{the express delivery service industry in China. You are to write a response about its recent development and its} \textit{impact on people’s lives. You will have 30 minutes to write the essay. You should write at least }\textit{\uline{120 }}\textit{words but no} \textit{more than }\textit{\uline{180 }}\textit{words.}


\subsection*{\textbf{Part II} \textbf{Listening Comprehension} \textbf{(25 minutes)}}


\subsection*{\textbf{Section A}}


\textbf{Directions:} \textit{In this section, you will hear three news reports. At the end of each news report, you will hear two or} \textit{three questions. Both the news report and the questions will be spoken only once. After you hear a question, you} \textit{must choose the best answer from the four choices marked A), B), C) and D). Then mark the corresponding letter} \textit{on }\textbf{\textit{Answer Sheet 1}}\textit{ with a single line through the centre.}


\textbf{Questions 1 and 2 are based on the news report you have just heard.}


1.


\begin{itemize}[leftmargin=2em,itemsep=.2em,topsep=.4em]\item[A.] She is making a special effort to be recognized.
\item[B.] She has been living in Fayetteville for 25 years.
\item[C.] She has been driving a school bus ever since 2009.
\item[D.] She is the longest-serving bus driver in Fayetteville.
\end{itemize}

2.


\begin{itemize}[leftmargin=2em,itemsep=.2em,topsep=.4em]\item[A.] The importance of their service to the city.
\item[B.] The number of riders they serve in the city.
\item[C.] The new proclamation issued early this week.
\item[D.] The significance of the event to take place March 18.
\end{itemize}

\textbf{Questions 3 and 4 are based on the news report you have just heard.}


3.


\begin{itemize}[leftmargin=2em,itemsep=.2em,topsep=.4em]\item[A.] Budgeted \$56.7 million for the coming school year.
\item[B.] Requested to raise teachers’ salaries by 3 percent.
\item[C.] Proposed a spending increase by 5.99 percent.
\item[D.] Kept raising funds for at least 8 budget cycles.
\end{itemize}

4.


\begin{itemize}[leftmargin=2em,itemsep=.2em,topsep=.4em]\item[A.] Decrease in salaries for school administrators.
\item[B.] Reduction in federal and state funding.
\item[C.] Firing of less qualified faculty members.
\item[D.] Closing down of some less competitive schools.
\end{itemize}

\textbf{Questions 5 to 7 are based on the news report you have just heard.}


5.


\begin{itemize}[leftmargin=2em,itemsep=.2em,topsep=.4em]\item[A.] They often take place in the Scottish highlands.
\item[B.] They usually happen in the middle of the night.
\item[C.] Their damage can be as serious as that of bombs.
\item[D.] Their loud claps can be confused with explosions.
\end{itemize}

6.


\begin{itemize}[leftmargin=2em,itemsep=.2em,topsep=.4em]\item[A.] They sent out a message to calm them down.
\item[B.] They promised to stop the construction work.
\item[C.] They told them the military exercise was over.
\item[D.] They set off to inspect the site of the plane crash.
\end{itemize}

\noindent\begin{tabularx}{\linewidth}{@{}lXXXX@{}}7. & \textbf{A.} Keep warm. & \textbf{B.} Take caution. & \textbf{C.} Avoid driving. & \textbf{D.} Use rail service.\\\end{tabularx}\par

\subsection*{\textbf{Section B}}


\textbf{Directions:} \textit{In this section, you will hear two long conversations. At the end of each conversation, you will hear} \textit{four questions. Both the conversation and the questions will be spoken only once. After you hear a question, you} \textit{must choose the best answer from the four choices marked A), B), C) and D). Then mark the corresponding letter} \textit{on }\textbf{\textit{Answer Sheet 1}}\textit{ with a single line through the centre.}


\textbf{Questions 8 to 11 are based on the conversation you have just heard.}


8.


\begin{itemize}[leftmargin=2em,itemsep=.2em,topsep=.4em]\item[A.] Why the man is so anxious to buy a new phone.
\item[B.] How the man could possibly afford a new phone.
\item[C.] How the man got into the habit of being wasteful.
\item[D.] Why the man is so fond of using his credit card.
\end{itemize}

9.


\begin{itemize}[leftmargin=2em,itemsep=.2em,topsep=.4em]\item[A.] It is not in fashion any more.
\item[B.] It goes out of order frequently.
\item[C.] It is not as fancy as his colleagues’.
\item[D.] It lacks functions office work requires.
\end{itemize}

10.


\begin{itemize}[leftmargin=2em,itemsep=.2em,topsep=.4em]\item[A.] It enables him to buy a new phone at a more favorable price.
\item[B.] It is a special offer to recent college graduates working in offices.
\item[C.] It is available to office workers who have a good enough credit score.
\item[D.] It allows him to borrow money without paying interest for six months.
\end{itemize}

11.


\begin{itemize}[leftmargin=2em,itemsep=.2em,topsep=.4em]\item[A.] She borrowed money to pay her debt.
\item[B.] She developed a habit of overspending.
\item[C.] She enjoys buying a lot of fancy things.
\item[D.] She regrets ignoring the woman’s advice.
\end{itemize}

\textbf{Questions 12 to 15 are based on the conversation you have just heard.}


12.


\begin{itemize}[leftmargin=2em,itemsep=.2em,topsep=.4em]\item[A.] Blog writing and natural resources saving.
\item[B.] Joint efforts to improve the local community.
\item[C.] Audience participation in \textit{The Morning Show}.
\item[D.] Home ownership and environmental protection.
\end{itemize}

13.


\begin{itemize}[leftmargin=2em,itemsep=.2em,topsep=.4em]\item[A.] Not a realistic size for families.
\item[B.] Not welcomed by most Americans.
\item[C.] Only fit for families of four or less.
\item[D.] Only suitable for renting to the poor.
\end{itemize}

14.


\begin{itemize}[leftmargin=2em,itemsep=.2em,topsep=.4em]\item[A.] It should be changed.
\item[B.] It isn’t easy to realize.
\item[C.] It will come true sooner or later.
\item[D.] It doesn’t appeal to average families.
\end{itemize}

15.


\begin{itemize}[leftmargin=2em,itemsep=.2em,topsep=.4em]\item[A.] They help to lower housing prices.
\item[B.] They help to save up for larger houses.
\item[C.] They are a good choice for many people.
\item[D.] They are vital to reducing waste worldwide.
\end{itemize}

\subsection*{\textbf{Section C}}


\textbf{Directions:}\textit{ In this section, you will hear three passages. At the end of each passage, you will hear three or four} \textit{questions. Both the passage and the questions will be spoken only once. After you hear a question, you must} \textit{choose the best answer from the four choices marked A), B), C) and D). Then mark the corresponding letter on} \textbf{\textit{Answer Sheet 1}}\textit{ with a single line through the centre.}


\textbf{Questions 16 to 18 are based on the passage you have just heard.}


16.


\begin{itemize}[leftmargin=2em,itemsep=.2em,topsep=.4em]\item[A.] Play outdoors and enjoy themselves.
\item[B.] Beat challenges and take risks.
\item[C.] Take care of their well-being.
\item[D.] Grow up strong and healthy.
\end{itemize}

17.


\begin{itemize}[leftmargin=2em,itemsep=.2em,topsep=.4em]\item[A.] Enable them to develop their motor skills.
\item[B.] Help them to conceal their frustration.
\item[C.] Reduce their stress and depression.
\item[D.] Prevent them from feeling lonely.
\end{itemize}

18.


\begin{itemize}[leftmargin=2em,itemsep=.2em,topsep=.4em]\item[A.] It begins early.
\item[B.] It proves demanding.
\item[C.] It calls for lifelong efforts.
\item[D.] It lies in nature experiences.
\end{itemize}

\textbf{Questions 19 to 21 are based on the passage you have just heard.}


19.


\begin{itemize}[leftmargin=2em,itemsep=.2em,topsep=.4em]\item[A.] When companies share their information online.
\item[B.] When their job search criteria are met fully.
\item[C.] When companies embrace greater diversity.
\item[D.] When they find job postings visually attractive.
\end{itemize}

20.


\begin{itemize}[leftmargin=2em,itemsep=.2em,topsep=.4em]\item[A.] Emails in the normal format.
\item[B.] Major companies in one location.
\item[C.] All companies in the same sector.
\item[D.] Various aspects of a company’s workforce.
\end{itemize}

\noindent\begin{tabularx}{\linewidth}{@{}lXXXX@{}}21. & \textbf{A.} Job security. & \textbf{B.} Potential pay raises. & \textbf{C.} Minimum base salary. & \textbf{D.} Information about diversity.\\\end{tabularx}\par

\textbf{Questions 22 to 25 are based on the passage you have just heard.}


\noindent\begin{tabularx}{\linewidth}{@{}lXXXX@{}}22. & \textbf{A.} Positive. & \textbf{B.} Original. & \textbf{C.} Important. & \textbf{D.} Surprising.\\\end{tabularx}\par

23.


\begin{itemize}[leftmargin=2em,itemsep=.2em,topsep=.4em]\item[A.] Dogs match their behavior with kids as much as with adults.
\item[B.] Kids and adults alike find in dogs a source of attachment.
\item[C.] Kids’ attitude to dogs is influenced by adults.
\item[D.] Dogs can help kids in many ways.
\end{itemize}

24.


\begin{itemize}[leftmargin=2em,itemsep=.2em,topsep=.4em]\item[A.] Dogs pay much attention to the kids they live with.
\item[B.] Dogs view both adults and kids as social companions.
\item[C.] Kids differ from adults in their behavior towards dogs.
\item[D.] Kids need to be aware of the risks in playing with dogs.
\end{itemize}

\noindent\begin{tabularx}{\linewidth}{@{}lXXXX@{}}25. & \textbf{A.} Understanding dogs. & \textbf{B.} Training dogs. & \textbf{C.} Taking care of dogs. & \textbf{D.} Associating with dogs.\\\end{tabularx}\par

\subsection*{\textbf{Part III Reading Comprehension (40 minutes)}}


\subsection*{\textbf{Section A}}


\textbf{Directions:} \textit{In this section, there is a passage with ten blanks. You are required to select one word for each blank} \textit{from a list of choices given in a word bank following the passage. Read the passage through carefully before} \textit{making your choices. Each choice in the bank is identified by a letter. Please mark the corresponding letter for} \textit{each item on }\textbf{\textit{Answer Sheet 2}}\textit{ with a single line through the centre. You may not use any of the words in the bank} \textit{more than once.}


Scientists have known that depriving adult mice of vision can increase the sensitivity of individual neurons (神经元) in the part of the brain devoted to hearing. New research from biologists at the University of Maryland \uline{26} that sight deprivation also changes the way brain cells \uline{27} with one another, shifting the mice’s sensitivity to different frequencies.


“This study \uline{28} what we are learning about how manipulating vision can have a \uline{29} effect on the ability of an animal to hear long after the window for auditory (听觉的) learning was thought to have \uline{30 },” said Patrick Kanold, senior author of the study.


It was once thought that the sensory regions of the brain were not \uline{31} after a critical period in childhood. This is why children learn languages much more \uline{32} than adults. Kanold’s earlier research disproved this idea by showing that depriving adult mice of vision for a short period increased the sensitivity of individual neurons in the auditory cortex (皮质), which is devoted to hearing.


Young brains wire themselves according to the sounds they hear frequently, assigning areas of the auditory cortex to \uline{33} frequencies based on what they are used to hearing. The researchers found that, in adult mice, a week in the dark also changed the \uline{34} of space to different frequencies.


“We don’t know why we are seeing these patterns,” Kanold said. “We \uline{35} that it may have to do with what the mice are paying attention to while they are in the dark.”


\begin{itemize}[leftmargin=2em,itemsep=.2em,topsep=.4em]\item[A.] adaptable
\item[B.] closed
\item[C.] distribution
\item[D.] interact
\item[E.] narration
\item[F.] neutral
\item[G.] permanently
\item[H.] prescribes
\item[I.] readily
\item[J.] registered
\item[K.] reinforces
\item[L.] revealed
\item[M.] significant
\item[N.] specific
\item[O.] speculate
\end{itemize}

\subsection*{\textbf{Section B}}


\textbf{Directions:} \textit{In this section, you are going to read a passage with ten statements attached to it. Each statement} \textit{contains information given in one of the paragraphs. Identify the paragraph from which the information is derived.} \textit{You may choose a paragraph more than once. Each paragraph is marked with a letter. Answer the questions by} \textit{marking the corresponding letter on }\textbf{\textit{Answer Sheet 2}}\textit{.} \textbf{Why it’s wrong to look at work-life balance as an achievement}


A) Few topics have been so endlessly analysed as work-life balance. The quest to attain this mysterious state has dominated discussion around careers for years — especially for working parents. The concept is often presented as something to achieve, or a goal to reach. And once you’ve reached it, congratulations: you’ve made it; you’re a successful human being of the 21st century.


B) But the problem is that we often tell ourselves: “I’m going to put in eight hours’ worth of work, and then I’m going to put in eight hours’ worth of ‘me time’, which will include my family, my hobbies, my workout, my everything,” says Anat Lechner, clinical associate professor of management at New York University. “I don’t think it’s such a simple formula.”


C) And, according to new findings, it may not be. Some researchers are now encouraging us to stop thinking about work-life balance as an achievement that you either hit or don’t. Instead, they suggest it may be more of a lifelong process — a continuous, never-ending exercise that requires self-awareness and timely adjustments. Researchers Ioana Lupu and Mayra Ruiz-Castro argue that work-life balance is “a cycle, not an achievement”. In their 2020 study, the researchers interviewed nearly 80 employees at two London-based firms — an equal number of men and women between the ages of 30 and 50, all with at least one dependent child — who worked in middle or senior management roles.


D) Although it sounds like the respondents had a lot in common, here’s what separated them: about 30\% of the men and 50\% of the women reported resisting working long hours. The other respondents, meanwhile, all worked long hours because they thought that’s what successful professionals should do.


E) Lupu and Ruiz-Castro looked at those who rejected the long hours and they found that those workers actually had strikingly similar strategies for maintaining their work-life balance. They had a tendency to reflect and question assumptions in the name of self-awareness and regularly took steps to adjust the things standing in their way to work-life balance.


F) Lupu and Ruiz-Castro identified five steps that the respondents in the study who had better work-life balance used in their jobs. First, they paused and reconsidered beliefs such as “I’m a professional, so I should work, work, work”, and asked themselves questions like, “What’s currently causing me stress?” Second, after identifying the cause, they zeroed in on their resultant emotions. Did they feel angry, sad, energised? Third, they reprioritised, asking “Is working long hours really worth cutting back on family time?”, for example. Fourth, they considered their alternatives: is there anything at work that could be changed to accommodate these new priorities? And finally, they implemented changes, like asking their supervisor for greater flexibility, or deciding not to take on every project that comes their way.


G) This five-step process is something anyone can adopt. Going through the steps, and constantly checking in with yourself, can help you shift and adapt your professional life to something that will better harmonise with your personal one. “Awareness of your emotional state is essential in order to determine the changes you want to make in your work and in your life,” says Lupu.


H) New York University’s Lechner agrees that finding that balance is an ongoing pursuit. It’s not simply about dividing up the hours in your day between work, the gym, kids and chores. If the underlying emotional sources of stress are still there, then the time you actually spend at home may not be enjoyable. “We come home and even though physically we are there, mentally we still may be processing things that happened at work. We’re not present,” she says. What we call “work-life balance” is actually just a substitute to having a sense of fulfillment and contentment.


I) Of course, finding that balance probably shouldn’t be something you have to do by yourself. Research by Erin Kelly, professor of work and organisation studies at MIT, shows companies and managers can play a key role in creating a better environment for workers. For her book \textit{Overload: How Good Jobs Went Bad} \textit{and What to Do about It}, she and co-author Phyllis Moen split more than 1,000 employees at a Fortune 500 company into two groups, one that worked under a management redesign and one that continued working within the existing management structure.


J) Under the management redesign, many steps were taken to ensure better work-life balance and prevent burnout (精疲力竭). Managers were regularly reminded to explicitly support their employees. Workers were allowed to make changes, like cancelling 9 am meetings. All of this was done in the name of increasing job satisfaction and giving workers greater flexibility, and to assure workers that it was something management was committed to. Unsurprisingly, Kelly and Moen found that employees in the redesign group reported less stress and less burnout. They were less likely to quit their jobs; indeed, over the next four years, they were 40\% less likely to quit than those who kept working under the old policies.


K) “Work-life balance is understood to be an individual’s response, so people think ‘it’s up to me to manage the craziness of my work life’ ”, says Kelly. But organisations need to examine the demands they’re placing on employees. “The root problem is not how the two pieces of work and life come together. It’s that we have unrealistic expectations of what we’re asked to do on the work side.” If your workplace isn’t an environment where work-life balance is possible in the first place, any effort you attempt to make toward it on a personal level will be in vain.


L) That’s a conversation that appears to be gathering pace. The new prevalence of remote and flexible working models will likely all play important roles in how we balance our professional and personal lives. And if it seems like finding that perfect balance remains difficult to achieve, the experts say that keeping some perspective can help. For millions of people, work is about being able to put food on the table. Talking about work-life balance “is a very privileged conversation”, says Lechner. “If we’re reflecting, maybe we should also reflect on that.”


36. According to a management expert, work-life balance is not as simple as giving equal amounts of time to


work and personal life.


37. Research found that those who are given greater flexibility at work are less stressed and more likely to


stay in their jobs.


38. Workers who rejected working long hours tended to make regular adjustments in order to achieve work-


life balance.


39. Talking about work-life balance is said to be a privilege reserved for the better-off, not for those who


barely make a living.


40. Knowing one’s emotional state is of utmost importance in deciding what changes to make for a better


work-life balance.


41. More female professionals reported being reluctant to work overtime than their male counterparts.


42. Without organisational support any personal effort to maintain work-life balance will be unsuccessful.


43. The question of how to achieve work-life balance has long been the main subject of discussion among


workers.


44. You may not actually experience emotional wellbeing at home if you remain occupied with what


happened at work.


45. Some researchers suggest that work-life balance is not a goal to achieve, but a process for life to be


adjusted promptly.


\subsection*{\textbf{Section C}}


\textbf{Directions: }\textit{There are 2 passages in this section. Each passage is followed by some questions or unfinished} \textit{statements. For each of them there are four choices marked A), B), C) and D). You should decide on the best} \textit{choice and mark the corresponding letter on }\textbf{\textit{Answer Sheet 2}}\textit{ with a single line through the centre.}


\subsection*{\textbf{Passage One}}


\textbf{Questions 46 to 50 are based on the following passage.}


As a university student, I’ve come to realise just how little I know about money. I’ve come to the brutal realisation that Australia’s youth are being done a great disservice by not receiving any consistent financial education.


Diving headfirst into the crash course of starting university, I’ve quickly tried to get myself up to speed with the financial terms of the modern era to help manage my personal finances.


I’ve read some financial education books, done some online learning, and have spoken to my nearest and dearest for their pearls of financial wisdom.


There are undoubtedly hundreds of great resources out there for those wishing to improve their financial literacy, but while I was researching, I still kept wishing that I didn’t have to play catch-up with the money world. I felt that I was sailing the financial seas with no skills and no life jacket!


However, after talking to my friends who felt the same, I quickly realised that it wasn’t my ignorance but the lack of financial education in our schooling system that is leaving high schoolers seriously behind in the modern world of money.


Let’s compare a theoretical financial education subject to the standard compulsory English class. On average, English may not be the most popular subject, but it’s consistently on the schedule throughout high school, with all students graduating fluent in English.


A financial education subject should do just the same. It shouldn’t be just a one-day event but a course integrated throughout the whole of high school that would allow students to gradually expand their financial literacy, and would prevent the need for a ‘catch up’ phase once we’re out on our own after graduation.


In the same way that learning a language or new skills takes time, building financial skills requires practice and years to gradually accumulate bits and pieces of knowledge.


Giving young people the opportunity to become familiar with the world of money would provide them with a great advantage to enter adulthood with confidence and security so that they are able to manage their own money and look after themselves.


46. What has the author come to realise since entering university?


\begin{itemize}[leftmargin=2em,itemsep=.2em,topsep=.4em]\item[A.] He needs a crash course on financial terms.
\item[B.] He is very much lacking in financial literacy.
\item[C.] It requires consistent education to be financially independent.
\item[D.] It is unrealistic to give all Australian youth a financial education.
\end{itemize}

47. How did the author feel in today’s money world?


\begin{itemize}[leftmargin=2em,itemsep=.2em,topsep=.4em]\item[A.] Badly equipped to survive.
\item[B.] Ignorant of financial literature.
\item[C.] Barely capable of moving ahead.
\item[D.] Overwhelmed by the resources online.
\end{itemize}

48. What did the author realise after talking to his friends?


\begin{itemize}[leftmargin=2em,itemsep=.2em,topsep=.4em]\item[A.] They were as keen as he was on financial matters.
\item[B.] The schooling system was to blame for his trouble.
\item[C.] High schoolers knew nothing about the modern financial world.
\item[D.] Financial courses were as unpopular as compulsory English classes.
\end{itemize}

49. What is the author’s idea of a financial education course?


\begin{itemize}[leftmargin=2em,itemsep=.2em,topsep=.4em]\item[A.] It should foresee students’ needs after graduation.
\item[B.] It should provide students with some basic knowledge.
\item[C.] It should be taught the same way as English is taught.
\item[D.] It should be integrated into high school education.
\end{itemize}

50. What would financial literacy do to young people?


\begin{itemize}[leftmargin=2em,itemsep=.2em,topsep=.4em]\item[A.] Allow them to enter adulthood with financial security.
\item[B.] Enable them to look after themselves without worrying about money.
\item[C.] Render them confident and secure in terms of money management.
\item[D.] Help them become familiar with the world of money.
\end{itemize}

\subsection*{\textbf{Passage Two}}


\textbf{Questions 51 to 55 are based on the following passage.}


Chocolates save us from many things, especially emotional distress. They comfort us in times of trouble, calming down a racing heart by channelling happy calories inside us. We all have faith in chocolates to delight us in an instant!


Recently, chocolate lovers were heartbroken as scientists claimed that they can become extinct by 2050! But hey, we have some happy news for you. Scientists can still help save chocolates from dying out! If you are not aware as to why scientists made the statement about the death of this wonderful thing, let us tell you the facts. Chocolate trees, whose seeds are used to make chocolate, grow in the tropical plant world and require very specific weather conditions to prosper.


Now, fifty percent of the world’s cocoa (可可) beans come from two countries in West Africa: Côte d’Ivoire and Ghana. Scientists believe that both of these countries will experience a 3.8 \(^{\circ}\mathrm{F}\) temperature increase by 2050 due to global warming, endangering the cacao (可可树) farms in the rainforests. These farms will then have to be shifted to cooler mountainous areas, which are the natural habitat of wildlife. This will lead to some tough decisions: whether to grow chocolate or save wildlife.


Unfortunately, the crisis of global warming has already had a serious negative impact on cacao farms’


yields, leading to the prices of chocolates skyrocketing.


Scientists, however, are trying to find a long-term and eco-friendly solution to this problem! They are trying to modify the species with a gene-editing technology, which will transform the seedlings into a species that survives even in a drier and warmer climate.


According to a report by \textit{The Business Insider}, in the University of California’s new bio-sciences building, tiny green cacao seedlings are lined up in refrigerated greenhouses for a new experiment by using a technology called CRISPR. By manipulating the DNA of plants, this technology is already being used across the world to make plants tougher and cheaper. Similarly, in this unconventional experiment, scientists will make tiny, precise changes to the DNA of the seedlings to make the cocoa crops survive in warmer and drier climates.


51. What do people believe chocolates can do?


\begin{itemize}[leftmargin=2em,itemsep=.2em,topsep=.4em]\item[A.] Cheer them up instantly.
\item[B.] Create happy calories.
\item[C.] Conceal emotional distress.
\item[D.] Relieve them of heart trouble.
\end{itemize}

52. What was scientists’ recent assertion about chocolates?


\begin{itemize}[leftmargin=2em,itemsep=.2em,topsep=.4em]\item[A.] They could become a rare treat in the near future.
\item[B.] They could calm people down a bit in times of crisis.
\item[C.] They could prevent people from getting heartbroken.
\item[D.] They could become unavailable in less than 30 years.
\end{itemize}

53. What would happen if the cacao farms were shifted to cooler mountainous areas?


\begin{itemize}[leftmargin=2em,itemsep=.2em,topsep=.4em]\item[A.] The natural habitat of wildlife there would be ruined.
\item[B.] The cacao farmers would have a tough time to adapt.
\item[C.] The rainforests would be shrinking dramatically.
\item[D.] The quality of cocoa beans would suffer greatly.
\end{itemize}

54. What do we learn about the cacao farms in the crisis of global warming?


\begin{itemize}[leftmargin=2em,itemsep=.2em,topsep=.4em]\item[A.] They try to seek help from gene-editing scientists.
\item[B.] They decide to move to cooler mountainous areas.
\item[C.] They have suffered a lot due to a decrease in produce.
\item[D.] They have benefited by raising prices of cocoa beans.
\end{itemize}

55. What are scientists trying to do in the University of California’s new bio-sciences building?


\begin{itemize}[leftmargin=2em,itemsep=.2em,topsep=.4em]\item[A.] Build rows of refrigerated greenhouses for research on cacao seedlings.
\item[B.] Gene-edit cacao seedlings for them to withstand a drier, warmer climate.
\item[C.] Produce chocolates with the latest gene-editing technology.
\item[D.] Transplant the genes of tougher plants to cacao seedlings.
\end{itemize}

\subsection*{\textbf{Part IV Translation (30 minutes)}}


\textbf{Directions: }\textit{For this part, you are allowed 30 minutes to translate a passage from Chinese into English. You} \textit{should write your answer on }\textbf{\textit{Answer Sheet 2}}\textit{.}


中国政府十分重视环境保护。近年来，中国在减少空气、水和土壤污染上取得了显著成效。为了不断改善人们的生活环境，中国采取了一系列有效措施，包括大力发展清洁能源，改善公共交通，推广共享单车，实施垃圾分类。通过这些措施，中国的城市和农村正在绿起来、美起来。中国还积极参与国际合作，为全球环境保护做出了重要贡献。

\end{document}